\section{为什么要开这门课}

这门课的第一层动机是“研究者正在和底层技术逐渐脱节”。
更早的时候，研究者往往自己实现模型、训练模型，至少会下载一个模型并做 fine-tuning；而现在很多工作流已经变成了对一个闭源大模型直接做 prompt。
这种抽象当然提高了生产力，也确实释放了大量研究效率，但它也带来了一个问题：抽象是有泄漏的。你未必真的理解系统在底层怎么工作，哪些部分是稳健的，哪些部分只是表面上能跑。

讲师把这个问题说得很直接：如果研究要继续推进，尤其是要做基础研究，就不能只停留在“字符串输入，字符串输出”的黑箱层面。
很多真正有价值的问题，需要你把数据、系统、模型这些环节重新拆开，再做协同设计。

\begin{knowledgebox}{这门课的现实背景}
\begin{itemize}
    \item 研究范式已经从“自己训练模型”转向“调用大模型”。
    \item 这提高了效率，但也让很多人远离了真正的机制理解。
    \item 课程存在的理由，是让这种底层理解重新变得可学、可练、可验证。
\end{itemize}
\end{knowledgebox}

